% 公共宏包与自定义命令
\usepackage[margin=2.5cm]{geometry}
\usepackage{amsmath,amssymb}
\usepackage{graphicx}
\graphicspath{{figures/}}
\usepackage{booktabs}
\usepackage{tabularx}
% GB/T 7714—2015 顺序编码制（biblatex/biber 配套样式）。
\usepackage[backend=biber,style=gb7714-2015,sorting=none,
  gbtype=true,doi=false,url=false,isbn=false]{biblatex}
\addbibresource{ref.bib}
\usepackage[hidelinks]{hyperref}
\ctexset{
  section = {name = {,、}, number = \chinese{section}},
  subsection = {name = {（,）}, number = \chinese{subsection}}
}
\setcounter{tocdepth}{2}
\setcounter{secnumdepth}{2}
\newcommand{\TODO}[1]{\textbf{【TODO：#1】}}

% 学位论文拟定目录：局部编号、悬挂缩进与可跨页边框。
\usepackage{enumitem}
\usepackage[breakable]{tcolorbox}
\newlist{thesisoutline}{enumerate}{5}
\setlist[thesisoutline]{parsep=0pt,partopsep=0pt,align=left}
\setlist[thesisoutline,1]{
  label={\bfseries 第\chinese{thesisoutlinei}章},
  leftmargin=4.2em,labelwidth=3.7em,labelsep=0.5em,
  topsep=0pt,itemsep=7pt
}
\setlist[thesisoutline,2]{
  label={\chinese{thesisoutlineii}、},
  leftmargin=0pt,labelwidth=2.4em,labelsep=0.5em,
  topsep=4pt,itemsep=4pt
}
\setlist[thesisoutline,3]{
  label={（\chinese{thesisoutlineiii}）},
  leftmargin=1.6em,labelwidth=3.3em,labelsep=0.5em,
  topsep=2pt,itemsep=2pt
}
\setlist[thesisoutline,4]{
  label=\arabic*.,
  leftmargin=1.3em,labelwidth=1.1em,labelsep=0.5em,
  topsep=2pt,itemsep=1pt
}
\setlist[thesisoutline,5]{
  label=\arabic{thesisoutlineiv}.\arabic*,
  leftmargin=1.8em,labelwidth=1.6em,labelsep=0.5em,
  topsep=2pt,itemsep=1pt
}
\newtcolorbox{thesisframework}{
  breakable,sharp corners,boxrule=0.5pt,
  colback=white,colframe=black,colbacktitle=white,coltitle=black,
  title={论文具体框架目录},fonttitle=\bfseries,
  left=10pt,right=10pt,top=7pt,bottom=7pt,
  before skip=10pt,after skip=10pt,
  fontupper=\small\linespread{1.15}\selectfont
}
